\documentclass[11pt,a4paper]{article}
\usepackage{fontspec}
\setmainfont{Arial}
\setsansfont{Arial}
\usepackage[ngerman]{babel}
\usepackage[margin=20mm,headheight=14pt]{geometry}
\usepackage{graphicx,xcolor,longtable,array,booktabs,needspace}
\usepackage{amsmath}
\usepackage[unicode,hidelinks]{hyperref}
\usepackage{fancyhdr}
\definecolor{accent}{HTML}{17364A}
\definecolor{teal}{HTML}{087C83}
\pagestyle{fancy}\fancyhf{}
\fancyhead[L]{\small Dissertation | Fragenkatalog und Lernskript}
\fancyfoot[L]{\small Überarbeitete Fachbegriffe | 24.09.2026}
\fancyfoot[R]{\thepage}
\renewcommand{\headrulewidth}{0.3pt}
\setlength{\parindent}{0pt}\setlength{\parskip}{6pt}
\setlength{\emergencystretch}{4em}
\raggedright
\renewcommand{\contentsname}{Inhaltsverzeichnis}
\setcounter{tocdepth}{1}
\newcommand{\topic}[1]{\clearpage\section*{\color{accent}#1}\addcontentsline{toc}{section}{#1}}
\newcommand{\question}[1]{\Needspace{6\baselineskip}\par\vspace{6pt}{\large\bfseries\color{accent}#1}\par\nobreak}
\begin{document}
\begin{titlepage}
\vspace*{22mm}{\Huge\bfseries\color{accent}Dissertation\par}
\vspace{6mm}{\LARGE Fragenkatalog und Lernskript\par}
\vspace{12mm}{\large Vorbereitung für Florian Rzepka\par}
\vspace{16mm}
Einheitlicher LaTeX-Satz mit englischen Fachbegriffen und deutschen Erklärungen.

Enthalten sind die Fragen und Antworten, Rechenübungen, der Abbildungsatlas, die Tabellenprüfung, die Modellübersicht und ein Begriffsschlüssel.

Frage F075 wurde verständlicher formuliert. Die Modellübersicht und ihre offenen Nachweispunkte bleiben erhalten. Die Abbildungen zeigen unveränderte Originalausschnitte der Dissertation.

\vfill Stand der Überarbeitung: 24. September 2026.
\end{titlepage}
\tableofcontents
